\documentclass[11pt]{amsart}

\usepackage[T1]{fontenc}
\usepackage[utf8]{inputenc}
\usepackage{amsmath,amssymb,amsthm}
\usepackage[margin=1in]{geometry}
\usepackage{hyperref}

\newcommand{\OO}{\mathbb{O}}
\newcommand{\RR}{\mathbb{R}}
\newcommand{\ZZ}{\mathbb{Z}}
\newcommand{\LL}{\Lambda}
\newcommand{\Nrm}{N}

\title[Historical appendix (expanded): integral Cayley numbers 1923--1946]%
{Expanded draft of an appendix for\\
\emph{An order on the Leech lattice from a $\mathbb{Z}_3$-symmetric
triple-octonion product}\\
\large Historical note: integral Cayley numbers, 1923--1946}
\author{Claude Opus 4.7}
\address{Anthropic}
\thanks{Drafted by the AI agent Claude Opus 4.7 (Anthropic) at the
direction of Jens K\"oplinger, as part of the systematic investigation
recorded in this repository (\texttt{prompt\_logs/}).  The short version
that appears in the main paper (\texttt{paper/historical\_appendix.tex})
is owned by the principal investigator; this expanded form is a
standalone reading copy and is intentionally longer than the version
that will be merged into the manuscript.}
\date{22 May 2026}

\begin{document}
\maketitle

\begin{abstract}
This is an expanded draft of the historical appendix.  It is written
\emph{solely from the primary sources}: the four papers Dickson (1923),
Kirmse (1924), Mahler (1942), and Coxeter (1946), each obtained and
re-derived in exact arithmetic in this project
(\texttt{python\_project/src/verify\_dickson\_1923.py},
\texttt{verify\_kirmse\_1924.py},
\texttt{verify\_mahler\_1942.py},
\texttt{verify\_coxeter\_1946.py};
companion notes in \texttt{paper/reviews/}).  No secondary source is taken
on trust.  A short form of this appendix, targeted at the main paper's
length budget, is in \texttt{paper/historical\_appendix.tex}.
\end{abstract}

% ===================================================================
\section*{Setting}
% ===================================================================

The integral-Cayley-number construction -- a maximal order in the real
octonion algebra, supported on the $E_8$ lattice -- was developed in four
papers between 1923 and 1946.  Each contributes one piece.  Together they
fix three facts that underlie the present construction: there are
\emph{seven} maximal orders; each is obtained from any one by a
transposition of two imaginary basis units (a linear involution of~$\RR^8$);
and the underlying lattice is~$E_8$.  The third fact and the count are
already in Dickson; the second is the Coxeter--Bruck remedy of Kirmse.

% ===================================================================
\section*{Dickson 1923}
% ===================================================================

L.\,E.~Dickson \cite{Dickson1923}, in Section~19 of \emph{A new simple theory
of hypercomplex integers}, gave the first construction.  He built the
octonions by Cayley--Dickson doubling of the quaternions, $x = q + Qe$ with
$q, Q$ quaternions, and the product (his equation~(24))
\[
  (q + Qe)(r + Re) \;=\; (qr - \bar R Q) + (Rq + Q \bar r)\, e,
\]
and then determined the maximal sets of integers closed under his three
axioms -- closure C, contains unit~U, integral norm~N (his Section~6).
From the norm form he derived that integer coordinates are either integers
or half-integers, ran a case analysis, and exhibited (his~(29)) the basis
\[
  i,\; j,\; k,\; p,\; e,\; W,\; Z,\; V,
  \qquad p = \tfrac{1}{2}(1 + i + j + k),
\]
of an order, his \emph{system~(30)} (here denoted~$O_D$).  All $64$ products
of the basis~(29) lie in~$O_D$; the system is genuinely closed
\cite{Dickson1923}.

Dickson then stated as \textbf{Theorem~XV}:
\begin{quote}
\emph{The only maximal systems of integers of Cayley's algebra having
properties~C, U, N are~(30) and the two systems obtained from it by
cyclic permutations of~$i, j, k$.}
\end{quote}
This undercounts.  Re-derived in this project, the true number of maximal
orders for Dickson's product is~\emph{seven}, not three.  The cause is in
Dickson's own derivation: on p.\,321, having reached a case in which the
order contains the Hurwitz quaternions $\tfrac{1}{2}(\pm 1 \pm i \pm j \pm
k)$, he writes
\begin{quote}
\emph{``\ldots and hence all of Hurwitz's integral quaternions.  We shall
assume that the latter occur in our set in all cases.''}
\end{quote}
The assumption that every maximal order contains the Hurwitz quaternions is
stronger than C, U, N.  Of the seven maximal orders, \emph{three} contain
those quaternions, and those three are precisely $O_D$ and its two
$(i\,j\,k)$-cyclic images -- exactly Dickson's Theorem~XV triple.  The four
he missed are exactly the maximal orders that do not contain the integral
quaternions.

So Dickson's \emph{construction} is correct; his \emph{count} is too small
by four, by a qualifier that he stated openly in the derivation but that was
absorbed into the theorem.

% ===================================================================
\section*{Kirmse 1924}
% ===================================================================

J.~Kirmse \cite{Kirmse1924}, working in Germany the following year, took up
the problem with a different presentation.  He fixed an octonion
multiplication on the basis $\{1, i_1, \ldots, i_7\}$ by seven cyclic Fano
triples (his table~(1) on p.\,64) and sought the maximal sets of octonion
integers closed under multiplication.  He stated that there are
\emph{eight} such sets (\emph{``Man findet acht solche \ldots,''} p.\,70)
and described one in detail, his system $J_1$ (p.\,70):
\[
  J_1 \;=\; \big[\, J_0,\; a_1,\; a_2,\; a_3,\; a_4\,\big],
\]
where $J_0 = \ZZ\langle 1, i_1, \ldots, i_7\rangle$ and
\begin{align*}
  a_1 &= \tfrac{1}{2}(1 + i_1 + i_2 + i_3), &
  a_2 &= \tfrac{1}{2}(i_4 + i_5 + i_6 + i_7), \\
  a_3 &= \tfrac{1}{2}(1 + i_1 + i_4 + i_5), &
  a_4 &= \tfrac{1}{2}(1 + i_2 + i_4 + i_7).
\end{align*}

Kirmse's multiplication table is itself a correct octonion algebra: it is
norm-multiplicative and satisfies all the Cayley identities.  His error is
in the lattice~$J_1$: it is \emph{not closed} under his own multiplication.
The witness, computed in this project and matching Coxeter (1946) verbatim,
is
\[
  a_1 \cdot a_3
  \;=\; \tfrac{1}{2}(1 + i_1 + i_2 + i_3) \cdot
        \tfrac{1}{2}(1 + i_1 + i_4 + i_5)
  \;=\; \tfrac{1}{2}(i_1 + i_2 + i_4 + i_7),
\]
and the product is not among the 240 elements of $J_1$ of unit norm
(Kirmse p.\,76).  An enumeration of the $30$ candidate $E_8$-type lattices
under Kirmse's table yields exactly \emph{seven} closed ones; $J_1$ is the
spurious entry on Kirmse's list of eight.

So Kirmse's \emph{table} is correct; his \emph{exhibited order} is the
defective one, and his \emph{count} is too large by one.

% ===================================================================
\section*{Mahler 1942}
% ===================================================================

K.~Mahler \cite{Mahler1942}, eighteen years later, worked over a maximal
order (one of the seven, to be made explicit by Coxeter in 1946) and
developed its arithmetic.  His principal results are two.

\begin{itemize}
\item \textbf{Euclidean property} (his Theorem on p.\,127):
\emph{for every Cayley number~$X$ there is an integral Cayley number~$G$
with $\Nrm(X - G) \le \tfrac{1}{2}$.}  Mahler calls this ``not easy to
prove.''  It is sharper than Dickson's earlier $5/4$ in place
of~$\tfrac{1}{2}$.
\item \textbf{Principal ideals}: \emph{every left (or right) ideal in the
ring of integral Cayley numbers is principal, generated by a multiple (in
the ordinary sense) of an integral Cayley number of norm~$1$, $2$, or~$4$.}
\end{itemize}

Re-derived in this project, Mahler's results are correct as stated; the
printed paper contains one transcription error in equation~(10), without
consequence for the theorems.  Mahler cites Dickson~(1923) but not Kirmse;
that citation chain is how the priority of Dickson's construction reaches
Coxeter in the following paper.

% ===================================================================
\section*{Coxeter 1946}
% ===================================================================

H.\,S.\,M.~Coxeter \cite{Coxeter1946}, in \emph{Integral Cayley numbers},
identifies Kirmse's error and supplies a corrected order.  His Section~4
(``Kirmse's mistake'') reproduces Kirmse's table as his~(4.1) (the two
tables are identical, verified) and exhibits the same non-closure witness
as above:
\[
  \tfrac{1}{2}(1 + i_1 + i_2 + i_3) \cdot
  \tfrac{1}{2}(1 + i_1 + i_4 + i_5)
  \;=\; \tfrac{1}{2}(i_1 + i_2 + i_4 + i_7),
\]
which, Coxeter notes, ``is not one of the 240 units.''  Of the count he
writes: \emph{``Actually there are only seven, which presumably are the
remaining seven of his eight''} (p.\,566).

The remedy is due to \textbf{R.\,H.~Bruck}.  Coxeter writes:
\emph{``When I pointed this out to Bruck, he sent me a revised description
of such a domain.  Bruck's domain~$J$ can be derived from~$J_1$ by
transposing two of the~$i$'s.''}  That is, the correction is a transposition
of two imaginary basis units of the octonions -- a linear involution
of~$\RR^8$.  Any of the $\binom{7}{2} = 21$ such transpositions repairs
Kirmse's~$J_1$.  Coxeter further notes that the $21$ ``fall into $7$ sets of
$3$, each set having the same effect,'' yielding the seven maximal orders;
this combinatorial structure is verified in this project (each set of $3$
is three disjoint transpositions fixing one common imaginary unit, the
``special unit'' of the corresponding order).

Coxeter then develops the geometry of the corrected order.  Its $240$
norm-$1$ units are the vertices of Gosset's polytope~$4_{21}$; the order
itself is the eight-dimensional honeycomb~$5_{21}$, the densest sphere
packing in eight dimensions (his Sections~6--13).  He proves the corrected
order maximal in his Section~11.

A postscript (Section~14) records that Coxeter learned of Dickson and
Mahler only after the manuscript was written, when Olga Taussky Todd drew
his attention to Mahler, which in turn cites Dickson.  Coxeter there
identifies the cause of Dickson's Theorem~XV undercount independently:
\begin{quote}
\emph{``\ldots\ he asserted [9; 334] that there are only three systems which
include all the eight basic units \ldots\ He misused the remaining four of
the seven systems by adding the very artificial requirement that the
integral Cayley numbers should include the integral quaternions
$\tfrac{1}{2}(\pm 1 \pm i \pm j \pm k)$.''}
\end{quote}
And of Mahler's ``not easy to prove'' Euclidean theorem, Coxeter observes
that it is an immediate consequence of his own Section~11.  By 1946 the
construction is complete.

% ===================================================================
\section*{Synthesis}
% ===================================================================

The four papers together fix the construction the present paper rests on:

\begin{itemize}
\item \textbf{1923 (Dickson).}  An integral-Cayley-number maximal order
exists; the lattice is $E_8$; the count, in the form Dickson's theorem
gives, is three (under his stated extra assumption).
\item \textbf{1924 (Kirmse).}  An independent treatment; correct
multiplication; an enumerated count of eight that includes one non-closed
candidate.
\item \textbf{1942 (Mahler).}  Arithmetic of the order: Euclidean
with~$\Nrm(X-G) \le \tfrac{1}{2}$; principal ideals.
\item \textbf{1946 (Coxeter, with Bruck).}  Kirmse's $J_1$ is not closed
(an explicit one-line witness); Bruck's remedy is a transposition of two
imaginary basis units; the $21$ such transpositions yield exactly the seven
maximal orders; Coxeter's postscript supplies the missing connection to
Dickson and Mahler.
\end{itemize}

\noindent
The count converges over the four papers -- $3 \to 8 \to 7$ -- and the
correction that turns Kirmse's~$J_1$ into a genuine order is the
\emph{Coxeter--Bruck transposition} of two imaginary coordinate axes: a
linear involution of~$\RR^8$.

This transposition is the historical precedent for the present paper's
$\sigma$ (Definition of the transposition-twisted product, Section~3).  The
two are the same kind of map -- a coordinate transposition -- differing
only in what they are applied to.  Coxeter--Bruck is applied to a
maximal-order candidate that sits asymmetrically in~$\RR^8$, and so changes
which $E_8$ placement is considered.  The present~$\sigma$ is applied where
the order is~$L = D_8^+$, the coordinate-symmetric placement; hence
$\sigma(L) = L$ and the $E_8$-level closure is preserved automatically
($L \cdot_\sigma L = \sigma(L \cdot L) \subseteq L$).  The
\emph{Coxeter--Bruck transposition} thus appears below as the same operation
at a different stratum: at the Leech-lattice level rather than at the
maximal-order level, redirecting Wilson's sublattice conditions rather than
repairing the choice of order.

% ===================================================================
\begin{thebibliography}{9}
% ===================================================================

\bibitem[Dickson1923]{Dickson1923}
L.\,E.~Dickson,
\emph{A new simple theory of hypercomplex integers},
J.~Math.\ Pures Appl.\ (9) \textbf{2} (1923), 281--326.

\bibitem[Kirmse1924]{Kirmse1924}
J.~Kirmse,
\emph{\"Uber die Darstellbarkeit nat\"urlicher ganzer Zahlen als Summen von
acht Quadraten und \"uber ein mit diesem Problem zusammenh\"angendes
nichtkommutatives und nichtassoziatives Zahlensystem},
Ber.\ Verh.\ S\"achs.\ Akad.\ Wiss.\ Leipzig, Math.-Phys.\ Kl.\
\textbf{76} (1924), 63--82.

\bibitem[Mahler1942]{Mahler1942}
K.~Mahler,
\emph{On ideals in the Cayley--Dickson algebra},
Proc.\ Roy.\ Irish Acad.\ Sect.\ A \textbf{48} (1942), 123--133.

\bibitem[Coxeter1946]{Coxeter1946}
H.\,S.\,M.~Coxeter,
\emph{Integral Cayley numbers},
Duke Math.\ J.\ \textbf{13} (1946), 561--578.

\end{thebibliography}

\end{document}
